\documentclass[11pt]{article}

\usepackage[margin=1.05in]{geometry}
\usepackage{amsmath,amssymb,amsthm,mathtools}
\usepackage[T1]{fontenc}
\usepackage{lmodern}
\usepackage{microtype}
\usepackage[hidelinks]{hyperref}

\newtheorem*{theorem}{Theorem}
\newtheorem*{lemma}{Lemma}
\newtheorem*{remark}{Remark}

\newcommand{\cube}{\{0,1\}}
\newcommand{\R}{\mathbb{R}}
\newcommand{\Hstar}{H^{\ast}}
\newcommand{\sgn}{\operatorname{sgn}}
\newcommand{\im}{\operatorname{Im}}

\title{\vspace{-2.0em}Positive Parallel Cuts Give a Head Upper Bound}
\author{}
\date{}

\begin{document}
\maketitle
\vspace{-2.6em}

\section*{Setup}

Work on the Boolean cube $\cube^n$.  A Boolean function is a map
\[
f:\cube^n\to\cube.
\]
The quantity $\Hstar(f)$ is the minimum number of attention heads needed by a
single-layer attention-only model, with a final linear threshold readout, to
compute $f$ exactly on the cube.

For the present note, we only need the following scalar normal form for one
head.  Choose parameters
\[
\gamma>0,\qquad \rho_i>0,\qquad \alpha>0,\qquad \alpha\neq1,
\]
and choose arbitrary real parameters
\[
\eta,\quad m_1,\ldots,m_n,\quad \delta.
\]
After composing the head output with the final linear readout, a single head can
contribute the scalar atom
\[
\phi(x)=
\frac{\eta+\sum_{i=1}^{n}\rho_i\alpha^{x_i}(m_i+\delta x_i)}
     {\gamma+\sum_{i=1}^{n}\rho_i\alpha^{x_i}},
\]
where $x\in\cube^n$.

Let $H$ be a nonnegative integer.  To prove $\Hstar(f)\leq H$, it is enough to
choose $H$ atoms
\[
\phi_1,\ldots,\phi_H
\]
and a constant $c\in\R$ such that the real score
\[
S(x)=c+\sum_{h=1}^{H}\phi_h(x)
\]
satisfies
\[
f(x)=1 \iff S(x)>0
\qquad
\text{for every }x\in\cube^n.
\]

\section*{Positive parallel cuts}

A positive parallel statistic is an affine-linear statistic with strictly
positive variable coefficients:
\[
t(x)=\sum_{i=1}^{n}\lambda_i x_i,
\qquad
\lambda_i>0.
\]
Its level sets
\[
t(x)=a
\]
for real values $a$ are parallel affine hyperplanes in $\R^n$.  The positivity
assumption means that all these hyperplanes have normal vector in the positive
orthant.

Suppose there is a function
\[
F:\im(t)\to\cube
\]
such that $f$ factors through $t$:
\[
f(x)=F(t(x)).
\]
Write the distinct values of $t$ on the cube in increasing order.  Thus, for
some integer $M\geq1$, there are real numbers $\tau_0,\ldots,\tau_{M-1}$ with
\[
\im(t)=\{\tau_0<\tau_1<\cdots<\tau_{M-1}\}.
\]
For each $m\in\{0,\ldots,M-1\}$, define the sign
\[
\sigma_m=
\begin{cases}
+1,&F(\tau_m)=1,\\
-1,&F(\tau_m)=0.
\end{cases}
\]
The number of label changes along this one-dimensional profile is
\[
C_t(F):=
\left|\{m\in\{1,\ldots,M-1\}:\sigma_{m-1}\neq\sigma_m\}\right|.
\]

Let $\mathcal A_+(f)$ be the set of pairs $(t,F)$ satisfying all of the
conditions just stated:
\[
\mathcal A_+(f):=
\left\{
(t,F):
t(x)=\sum_{i=1}^{n}\lambda_i x_i,\ \lambda_i>0,\ F:\im(t)\to\cube,\ f=F\circ t
\right\}.
\]
For every pair $(t,F)\in\mathcal A_+(f)$, the number $C_t(F)$ is defined as
above.  Define the positive-parallel complexity by minimizing this count:
\[
P_{\parallel}^{+}(f):=
\min_{(t,F)\in\mathcal A_+(f)} C_t(F).
\]
This minimum is finite.  For example,
\[
t(x)=\sum_{i=1}^{n}2^{i-1}x_i
\]
is injective on $\cube^n$, so every $f$ factors through this positive statistic.

\begin{theorem}
For every Boolean function $f:\cube^n\to\cube$,
\[
\Hstar(f)\leq P_{\parallel}^{+}(f).
\]
\end{theorem}

\section*{Proof}

Fix a pair $(t,F)\in\mathcal A_+(f)$.  By definition of $\mathcal A_+(f)$,
there are positive coefficients $\lambda_1,\ldots,\lambda_n$ such that
\[
t(x)=\sum_{i=1}^{n}\lambda_i x_i,
\qquad
\lambda_i>0,
\]
and $f(x)=F(t(x))$ for every $x\in\cube^n$.

Let $M$ be the number of distinct values of $t$ on the cube.  Write these values
as
\[
\im(t)=\{\tau_0<\tau_1<\cdots<\tau_{M-1}\}.
\]
For each $m\in\{0,\ldots,M-1\}$, let
\[
\sigma_m=
\begin{cases}
+1,&F(\tau_m)=1,\\
-1,&F(\tau_m)=0.
\end{cases}
\]
Let
\[
C:=C_t(F).
\]
We prove
\[
\Hstar(f)\leq C.
\]
The theorem follows by minimizing over all positive-parallel representations.

If $C=0$, then $F$ is constant on $\im(t)$, so $f$ is constant on the cube.
Thus $\Hstar(f)=0$.  From now on assume $C\geq1$.

\subsection*{Step 1. A univariate sign polynomial}

Let
\[
J:=\{m\in\{1,\ldots,M-1\}:\sigma_{m-1}\neq\sigma_m\}.
\]
Thus $|J|=C$.  For every $m\in J$, choose a cutpoint
\[
a_m\in(\tau_{m-1},\tau_m).
\]
Define
\[
p(z):=\sigma_0\prod_{m\in J}(a_m-z).
\]
Then $p$ changes sign exactly once between $\tau_{m-1}$ and $\tau_m$ for each
$m\in J$, and it changes sign nowhere else between consecutive image values.
Hence
\[
\sgn p(\tau_m)=\sigma_m
\qquad
\text{for every }m=0,\ldots,M-1.
\]
Equivalently,
\[
F(\tau_m)=1
\iff
p(\tau_m)>0.
\]

\subsection*{Step 2. Put the sign polynomial over a positive denominator}

Choose distinct positive shifts
\[
r_1,\ldots,r_C>0
\]
and set
\[
B(z):=\prod_{j=1}^{C}(z+r_j).
\]
Since every $\tau_m$ is nonnegative,
\[
B(\tau_m)>0
\qquad
\text{for every }m.
\]
Therefore the rational function
\[
R(z):=\frac{p(z)}{B(z)}
\]
has the same signs as $p$ on the finite set $\im(t)$:
\[
\sgn R(\tau_m)=\sigma_m.
\]

Because $\deg p=\deg B=C$ and the roots $-r_1,\ldots,-r_C$ of $B$ are simple,
partial fractions give real constants
\[
c,b_1,\ldots,b_C\in\R
\]
such that
\[
R(z)=c+\sum_{j=1}^{C}\frac{b_j}{z+r_j}.
\]
Substituting $z=t(x)$ gives
\[
R(t(x))=
c+\sum_{j=1}^{C}\frac{b_j}{t(x)+r_j}.
\]
This score has the correct signs on the cube:
\[
f(x)=1
\iff
F(t(x))=1
\iff
R(t(x))>0.
\]

\subsection*{Step 3. Each reciprocal term is one head}

It remains to show that, for any fixed $r>0$ and $b\in\R$,
\[
x\longmapsto \frac{b}{t(x)+r}
\]
is a valid one-head atom.

Let
\[
\Lambda:=\sum_{i=1}^{n}\lambda_i.
\]
Choose
\[
\alpha>1+\frac{\Lambda}{r},
\qquad
\rho_i:=\frac{\lambda_i}{\alpha-1},
\qquad
\gamma:=r-\frac{\Lambda}{\alpha-1}.
\]
Then
\[
\alpha>1,\qquad \rho_i>0,\qquad \gamma>0.
\]
On the Boolean cube,
\[
\alpha^{x_i}=1+(\alpha-1)x_i.
\]
Therefore
\[
\begin{aligned}
\gamma+\sum_{i=1}^{n}\rho_i\alpha^{x_i}
&=
r-\frac{\Lambda}{\alpha-1}
+\sum_{i=1}^{n}\frac{\lambda_i}{\alpha-1}
\left(1+(\alpha-1)x_i\right)\\
&=
r+\sum_{i=1}^{n}\lambda_i x_i\\
&=
r+t(x).
\end{aligned}
\]
Now take the numerator of the one-head atom to be the constant $b$, by setting
\[
\eta=b,\qquad m_i=0,\qquad \delta=0.
\]
The atom becomes exactly
\[
\phi(x)=\frac{b}{r+t(x)}.
\]

Thus each summand
\[
\frac{b_j}{t(x)+r_j}
\]
costs one head.  The constant $c$ is absorbed into the final readout threshold.
Consequently the score
\[
R(t(x))=
c+\sum_{j=1}^{C}\frac{b_j}{t(x)+r_j}
\]
is realized with $C$ heads and computes $f$ after thresholding at zero.  Hence
\[
\Hstar(f)\leq C_t(F).
\]
Minimizing over all positive-parallel representations gives
\[
\Hstar(f)\leq P_{\parallel}^{+}(f).
\]
\qed

\section*{What the positivity assumption is doing}

The proof used $\lambda_i>0$ in exactly one place: it made the denominator
\[
t(x)+r=r+\sum_i\lambda_i x_i
\]
positive on the cube with all slopes of the same sign.  That is precisely the
kind of denominator produced by a single softmax attention head in the scalar
normal form above.

If one allows arbitrary parallel hyperplanes
\[
L(x)=a_0+\sum_i a_i x_i
\]
with mixed-sign coefficients $a_i$, the same partial-fraction proof would
produce terms of the form
\[
\frac{b}{L(x)+r}.
\]
Here $b\in\R$ is a numerator coefficient and $r$ is a shift chosen so that the
denominator avoids zero on the cube.
Those denominators may have mixed-sign slopes.  They are not automatically
one-head denominators.  Thus the theorem proved here is the positive-parallel
bound
\[
\Hstar(f)\leq P_{\parallel}^{+}(f),
\]
not an unrestricted mixed-sign parallel bound.

\end{document}
